\chapter{Pilot USB Middleware Optimisation and Benchmarking}
\label{chap:cascade-usb-pilot}

The Universal Serial Bus (USB) device stack served as the first experimental test of our
method. We selected Human Interface Device (HID), Communications Device Class (CDC), and
Mass Storage Class (MSC) scenarios because the internal workbook records footprint and
initialisation measurements for all three. The pilot had two objectives: optimise the STM32
reference before comparison and determine which evidence controls were needed for later
campaigns.

\section{Evidence base and validation rules}
\label{sec:cascade-pilot-evidence}

The main Renesas tables use the clearly labelled \emph{Summary} sheet in the supplied pilot
workbook. The \emph{Renesas} sheet provides detailed cross-checks, and
\emph{Initialization Execution time} provides cycle-level traceability for STM32. For the
NXP study, slides~23--25 of the supplied \emph{ST Competitors MW Benchmark} presentation
are the authoritative record because they show the reviewed result tables and scenario
flowcharts. Where those images conflict with workbook cells, we retain the values visible
in the presentation.

The boards for this independent pilot campaign were:
\begin{itemize}
\item Renesas pilot: RA6M5 and STM32H573I-DK;
\item NXP pilot: MCX-N9 and STM32U575.
\end{itemize}
They were selected within the pilot work and are the only platforms to which the results in
this chapter apply.

\section{STM32 baseline-to-optimised results}
\label{sec:cascade-stm32-optimisation}

Table~\ref{tab:cascade-stm32-optimisation} separates the original STM32 USB projects from
their low-footprint configurations. This step prevents avoidable reference-project overhead
from being attributed to the competing platform.

\begin{table}[H]
\centering
\small
\caption{STM32 USB baseline and optimised results from the workbook Summary sheet.}
\label{tab:cascade-stm32-optimisation}
\begin{tabular}{lrrrrrr}
\toprule
& \multicolumn{2}{c}{ROM (bytes)} & \multicolumn{2}{c}{RAM (bytes)} & \multicolumn{2}{c}{Initialisation (ms)} \\
\cmidrule(lr){2-3}\cmidrule(lr){4-5}\cmidrule(lr){6-7}
Scenario & Base & Opt. & Base & Opt. & Base & Opt. \\
\midrule
HID & 54,095 & 52,655 & 43,733 & 36,657 & 0.97 & 0.71 \\
CDC & 33,260 & 33,284 & 18,830 & 17,422 & 0.60 & 0.50 \\
MSC & 70,480 & 70,476 & 108,109 & 69,197 & 2.77 & 2.24 \\
\bottomrule
\end{tabular}
\end{table}

The optimisation does not produce a uniform ROM reduction. CDC increases by 24~bytes and
MSC changes by only 4~bytes. Its stronger effect is visible in RAM and initialisation time,
especially for MSC, whose reported RAM falls from 108,109 to 69,197~bytes. The detailed
initialisation sheet records the corresponding STM32 cycle totals and permits the rounded
times in Table~\ref{tab:cascade-stm32-optimisation} to be traced to the captures.

\section{Optimised STM32 versus Renesas RA6M5}
\label{sec:cascade-pilot-renesas}

\begin{table}[H]
\centering
\small
\caption{RA6M5 and optimised STM32H573I-DK USB results from the workbook Summary sheet.}
\label{tab:cascade-renesas-results}
\begin{tabular}{lrrrrrr}
\toprule
& \multicolumn{2}{c}{ROM (bytes)} & \multicolumn{2}{c}{RAM (bytes)} & \multicolumn{2}{c}{Initialisation (ms)} \\
\cmidrule(lr){2-3}\cmidrule(lr){4-5}\cmidrule(lr){6-7}
Scenario & RA6M5 & STM32 & RA6M5 & STM32 & RA6M5 & STM32 \\
\midrule
HID & 61,380 & 52,655 & 77,896 & 36,657 & 1.59 & 0.71 \\
CDC & 45,224 & 33,284 & 53,090 & 17,422 & 10.98 & 0.50 \\
MSC & 132,712 & 70,476 & 121,252 & 69,197 & 3.84 & 2.24 \\
\bottomrule
\end{tabular}
\end{table}

For these three implementations, the optimised STM32 project has the lower reported ROM,
RAM, and initialisation value in every summary row. The largest timing separation appears
in CDC. The table does not isolate one cause: board choice, generated startup code,
middleware configuration, compiler output, and optimisation decisions remain part of the
measured implementations. The result therefore applies to the recorded RA6M5 and
STM32H573I-DK projects rather than to all Renesas and ST devices.

\section{STM32U575 versus NXP MCX-N9}
\label{sec:cascade-pilot-nxp}

The NXP summary reports two footprint boundaries. Table~\ref{tab:cascade-nxp-with-rcc}
includes the STM32 Reset and Clock Control (RCC) contribution, while
Table~\ref{tab:cascade-nxp-without-rcc} removes it. Presenting both prevents clock setup
from being hidden inside a nominal USB-stack comparison.

\begin{table}[H]
\centering
\small
\caption{STM32U575 and MCX-N9 footprint with STM32 RCC included, as shown in the reviewed presentation.}
\label{tab:cascade-nxp-with-rcc}
\begin{tabular}{lrrrr}
\toprule
& \multicolumn{2}{c}{ROM (bytes)} & \multicolumn{2}{c}{RAM (bytes)} \\
\cmidrule(lr){2-3}\cmidrule(lr){4-5}
Scenario & STM32 & MCX-N9 & STM32 & MCX-N9 \\
\midrule
HID & 16,377 & 13,728 & 2,422 & 1,702 \\
CDC & 17,013 & 13,549 & 3,065 & 1,360 \\
MSC & 35,025 & 26,177 & 13,556 & 11,859 \\
\bottomrule
\end{tabular}
\end{table}

\begin{table}[H]
\centering
\small
\caption{STM32U575 and MCX-N9 footprint with STM32 RCC excluded, as shown in the reviewed presentation.}
\label{tab:cascade-nxp-without-rcc}
\begin{tabular}{lrrrr}
\toprule
& \multicolumn{2}{c}{ROM (bytes)} & \multicolumn{2}{c}{RAM (bytes)} \\
\cmidrule(lr){2-3}\cmidrule(lr){4-5}
Scenario & STM32 & MCX-N9 & STM32 & MCX-N9 \\
\midrule
HID & 11,817 & 12,282 & 2,422 & 1,590 \\
CDC & 12,453 & 12,087 & 3,065 & 1,352 \\
MSC & 27,841 & 24,584 & 13,556 & 11,747 \\
\bottomrule
\end{tabular}
\end{table}

The Summary sheet lists an average STM32 ROM gap of $+26.22\,\%$ with RCC and
$+4.16\,\%$ without RCC. We preserve those reported aggregates rather than infer an
unstated averaging rule. Removing RCC changes the HID ROM ordering, while MCX-N9 retains
the lower reported RAM in all rows. The recorded STM32 RCC ROM contributions are 4,560,
4,560, and 7,184~bytes for HID, CDC, and MSC. The corresponding MCX-N9 clock/configuration
entries are 1,446/112, 1,462/8, and 1,593/112~bytes of ROM/RAM.

\subsection{Initialisation measurements}

\begin{table}[H]
\centering
\small
\caption{STM32U575 and MCX-N9 initialisation times shown in the reviewed presentation.}
\label{tab:cascade-nxp-timing}
\begin{tabular}{lrr}
\toprule
Scenario & STM32U575 (ms) & MCX-N9 (ms) \\
\midrule
HID device & 10 & 5 \\
CDC device & 9.993 & 8.159 \\
MSC host & 112 & 100.03 \\
\bottomrule
\end{tabular}
\end{table}

MCX-N9 has the lower displayed initialisation time in each scenario. The presentation
resolves the workbook's ambiguous MSC label: slide~24 identifies the case as a bare-metal
USB MSC host and shows the measured initialisation path through clock configuration, USB
configuration, and FATFS initialisation before the flash-disk insertion check. We therefore
report it as MSC host initialisation rather than as \texttt{MX\_USB\_Device\_Init()}.

\subsection{Resolution of workbook conflicts}

The workbook contains conflicting alternatives for three footprint entries. The reviewed
presentation resolves them: slide~23 shows 16,377~bytes for STM32 HID ROM and
13,549~bytes for MCX-N9 CDC ROM, while slide~24 shows MSC RAM of 13,556~bytes for
STM32U575 and 11,859~bytes for MCX-N9 with RCC included. The same slide shows
11,747~bytes for MCX-N9 after its clock-configuration RAM is removed. These displayed
values are used in Tables~\ref{tab:cascade-nxp-with-rcc} and
\ref{tab:cascade-nxp-without-rcc}; the contradictory workbook cells are not used.

\section{Pilot synthesis and tool motivation}
\label{sec:cascade-pilot-synthesis}

The pilot established three controls for the later work. We must optimise and document the
reference before comparing it, report the software boundary explicitly, and preserve
conflicts between source artifacts. It also exposed how much repeated effort was required
to classify linker-map content, align projects, and track provenance. Those limitations
became the requirements for MCU Workflow Studio in Chapter~\ref{chap:footprint-tool} and
the MCU Benchmark Copilot Agent in Chapter~\ref{chap:benchmark-agent}.